% Generated draft: requires mathematical and source review.
\documentclass[11pt]{article}
\usepackage[margin=25mm]{geometry}
\usepackage{fontspec}
\setmainfont{Latin Modern Roman}
\usepackage{amsmath,amssymb,mathtools,unicode-math}
\setmathfont{Latin Modern Math}
\usepackage{enumitem,needspace,xurl,hyperref}
\setlength{\parindent}{0pt}
\setlength{\parskip}{5pt}
\setlength{\emergencystretch}{3em}
\newcommand{\N}{\mathbb{N}}
\newcommand{\Z}{\mathbb{Z}}
\newcommand{\Q}{\mathbb{Q}}
\newcommand{\R}{\mathbb{R}}
\newcommand{\C}{\mathbb{C}}
\newcommand{\F}{\mathbb{F}}
\newcommand{\A}{\mathbb{A}}
\newcommand{\PP}{\mathbb{P}}
\newcommand{\E}{\mathbb{E}}
\newcommand{\eps}{\varepsilon}
\newcommand{\dd}{\,\mathrm{d}}
\begin{document}
\section*{UnsolvedMath ID: 20003119}
\textbf{Source problem code:} catalogue entry
\subsection{Definitions and mathematical statement}
Identify the marked complete finite-area once-punctured-torus Teichm\"uller space \(\mathcal T(S_{1,1})\) with the upper half-plane as a parameter space. The question asks for an explicit description of its \emph{hyperbolic-surface} Thurston metric, not the Poincar\'e metric of that half-plane or a metric on flat tori.

The supplied aid uses positive Fricke coordinates \((x,y,z)\), the trace coordinates of a fixed intersecting curve basis and its Farey neighbor, on
\[
x^2+y^2+z^2=xyz,\qquad x,y,z>2.
\]
Primitive slopes \(r\in\Q\cup\{\infty\}\), with primitive integer pairs identified up to sign, index essential simple closed curves. Let \(T_r\) be their positive trace functions, initialized by \(x,y,z\) on a Farey triangle. For adjacent triangles sharing the two slopes with functions \(T_a,T_b\), the opposite functions satisfy
\[
T_c+T_{c'}=T_aT_b.
\]
The reported Markoff-tree formulas are
\[
\ell_r(X)=2\operatorname{arcosh}(T_r(X)/2),\qquad
 d_{\rm Th}(X,Y)=\log\sup_r
\frac{\operatorname{arcosh}(T_r(Y)/2)}{\operatorname{arcosh}(T_r(X)/2)}.
\]
This is an exact supremum over all primitive slopes, not a claim that finitely many recursive steps always certify its value. Matching Fricke and upper-half-plane markings is necessary; the coordinate dictionary is not supplied here.

For a tangent vector \(v\) to the Fricke surface the corresponding reported infinitesimal formula is
\[
\|v\|_{\rm Th,X}=\sup_r
\frac{(\dd T_r)_X(v)}{\sqrt{T_r(X)^2-4}\,\operatorname{arcosh}(T_r(X)/2)}.
\]
The derivative covectors are obtained from the same recursion,
\[
\dd T_{c'}=T_b\,\dd T_a+T_a\,\dd T_b-\dd T_c.
\]
No absolute value is inserted in the asymmetric infinitesimal norm.

The aid gives an equivalent stable-length description. Let \(\|\cdot\|_X\) be the McShane--Rivin length norm on \(H_1(S_{1,1};\R)\), normalized so that a primitive integral class has the length of its simple geodesic. Then
\[
e^{d_{\rm Th}(X,Y)}=\sup_{v\ne0}\frac{\|v\|_Y}{\|v\|_X},
\]
the directed operator norm of the identity from the \(X\)-normed plane to the \(Y\)-normed plane. It also reports known Fenchel--Nielsen formulas for the two stretch paths associated to the two complete-lamination choices for a simple closed curve, giving exact unit-speed directed geodesics. Their explicit twist formulas and conventions are not supplied. These descriptions show why the asymmetric Thurston metric is different from the symmetric Poincar\'e metric on the same parameter space, without claiming a finite closed formula for the maximizing slope in every pair.
\subsection{Short English statement}
Describe the punctured-torus Thurston metric through recursively generated trace lengths, and its tangent norm through the derivatives of the same traces. The supplied stable-length formulation is a directed operator norm; it is distinct from both Poincar\'e geometry and the cited flat-torus setting.
\subsection{Sources}
\begin{itemize}
\item Problem 5.1, AIM workshop \emph{Lipschitz metric on Teichmueller space}, \url{https://aimath.org/pastworkshops/lipschitzteichproblems.pdf}. The source mentions Belkhirat, Papadopoulos and Troyanov's study of flat-torus Thurston metrics, without a full reference.
\item The research summary names McShane--Rivin stable length norms and known Fenchel--Nielsen stretch-path formulas; full titles, URLs and normalization details are not supplied.
\item Local UnsolvedMath registry, permanent ID 20003119, exact code \texttt{catalogue entry}; canonical location \texttt{aim-topology-notes.json}, notes[206], dated 14 August 2026; statement labeled exact.
\item Research attempt 1: partial result, problem partially solved at the run, rigor label proved, novelty confidence low; registry status partially solved. The Markoff-recursion and norm package is supplied as a machine-generated research aid, not peer reviewed; independent verification is required before citation.
\end{itemize}
\end{document}
